\chapter{Conclusion}

This project successfully demonstrated the fundamental mechanics of a buffer overflow vulnerability using the RISC-V architecture. By constructing a targeted assembly program within the Ripes simulator, we observed the complete exploit chain: from the initial unauthorized \texttt{sw} instruction, through memory corruption, \texttt{ra} overwrite, and ultimately, \texttt{PC} hijack. 

Visual verification at each pipeline stage conclusively confirmed that corrupting the return architectural slot effectively overrides processor control flow. When the illicit \texttt{jr ra} instruction executed, the pipeline encountered an indirect branch resolution delay, causing a measurable 2-cycle control hazard penalty as it speculatively processed unintended paths. This study bridges the theoretical understanding of control hazards with practical cybersecurity exploits, underlining the critical reality that processor architectures inherently execute instructions based on register and memory states without contextual validation. Consequently, advanced mitigation mechanisms involving both software compilation checks and hardware execution restrictions remain strictly necessary to ensure systematic memory integrity.
